% !PRINTED 125
\weberpage{125}

It is now further easy to see that $\gamma=g$ must be the case; for since $g$ is the \gesperrt{lower bound} of all values of $X$ within $(R)$, $\gamma$ cannot first of all be smaller than $g$.

But $\gamma$ also cannot be greater than $g$; for all values which $X$ assumes on the segment $MM'$ are $\ge\gamma$. If, however, $\gamma>g$, then one can choose a quantity $g'$ such that $\gamma>g'>g$, and, by the continuity of $X$ proved in \S\,31, VI, numbers $\alpha,\beta$ can be so determined that
\[
 \alpha<\eta<\beta
\]
and such that in the whole region $(P,\beta)-(P,\alpha)=(\alpha,\beta)$ ($QNN'Q'$ in Figure 2) $X$ remains greater than $g'$. Consequently the \gesperrt{lower bound} of $X$ in $(\alpha,\beta)$ is greater than $g$. Since now the \gesperrt{lower bound} of $X$ in $(P,\alpha)$ is also greater than $g$, it is also greater than $g$ in $(P,\beta)$, which is composed of $(P,\alpha)$ and $(\alpha,\beta)$; but $\beta$ belongs to $\mathfrak B$, and this would give a contradiction with the definition of $\mathfrak B$.

Thus there remains only the possibility that $\gamma=g$, and proposition 4 is thereby proved; namely, there is a value $\xi$ in the interior of $(R)$ for which
\[
 |f(\xi)|=g
\]
holds, so that $|f(x)|$ attains a minimum value. We now prove:

5. \gesperrt{If $\alpha$ is any value of $x$ for which $f(\alpha)$ is different from zero, then $h$ can be chosen so that}
\begin{equation}
 |f(\alpha+h)|<|f(\alpha)|
\tag{10}
\end{equation}
\gesperrt{results.}

It then follows that, if $f(\xi)$ were different from zero, $g$ could not be the minimum of the function $|f(x)|$; but since this has been proved, one must have
\begin{equation}
 f(\xi)=0,
\tag{11}
\end{equation}
and $\xi$ is a root of the equation $f(x)=0$. The fundamental theorem will therefore be proved.

In the proof of 5. we make use of the theorem already proved (\S\,32), that a pure equation always has a root.

Some of the derived functions $f'(\alpha),f''(\alpha),\ldots$ may vanish; the last, $f^{(n)}(\alpha)$, which is equal to $\Pi(n)a_0$, is however
%CONT>>
